% !TEX root = ../main.tex

\chapter{Introduction}


In this chapter, you must provide an overview of your thesis, written such that a non-engineer person (one of your parents) should be able to read and understand at high level.
\begin{itemize}

\item what is you main contribution? 
\item which methodology you used? (theory, simulation, experiments, literature overview)?
\item how the thesis is organized? which topic is covered in each chapter?

 \end{itemize}


\section{Motivation and Context}

Artificial Intelligence (AI) has become increasingly important in modern
computing systems worldwide, ranging from everyday applications such as
conversational assistants and recommendation systems to more complex tasks
involving prediction, classification, perception, autonomous decision-making,
and many others. Many of these applications rely on machine learning models
whose execution requires a large number of numerical operations. In particular,
Deep Neural Networks (DNNs) and other AI models are largely based on repeated
dot products and matrix multiply-accumulate operations, which represent
fundamental computational kernels in several AI workloads.

For this reason, the design of hardware architectures capable of accelerating
these mathematical operations has become a relevant topic in both academic and
industrial research. Specialized hardware accelerators can improve the execution
of matrix-based computations by exploiting parallelism and by using internal
arithmetic units optimized for specific numerical formats. This is especially
important because the efficiency of AI execution does not depend solely on the
algorithm itself, but also on how the underlying hardware represents, moves, and
processes the numerical data.

Therefore, the choice of the numerical format used during the execution of a
given AI model directly affects several important design aspects of the
associated accelerator circuit, including performance, circuit area, energy
efficiency, numerical accuracy, and resilience to faults. These aspects motivate
the study of hardware architectures capable of supporting and comparing multiple
numerical formats within the same computational framework.


\section{Problem Statement}

As mentioned previously, the problem addressed in this thesis concerns the
efficient design of dedicated accelerators for the numerous matrix
multiply-accumulate (MMA) operations involved in AI and DNN workloads.
Moreover, such calculations are performed by accelerators through the encoding
of the matrix elements into a specific numerical format. Depending on the
selected format, the hardware architecture of the accelerator determines how
data is computed, stored, moved, and approximated. Therefore, the choice of
the numerical format used during the execution of a given workload affects
both the hardware design of the accelerator and its related metrics.

At present, conventional standard formats such as half-precision and
single-precision floating-point, as well as integer representations, are
widely used in today's accelerators embedded in Graphics Processing Units
(GPUs), Central Processing Units (CPUs), and other hardware platforms. At the
same time, emerging number formats such as posit, fixed-point, logarithmic,
and reduced-size floating-point and integer formats also exist. However, their
usefulness is not always obvious until appropriate accelerators involving
their use are implemented, integrated, and tested.

Consequently, the problem grows in scope when the availability of multiple
possible numerical formats is taken into account. This thesis addresses this
problem through the design, integration, and testing, using hardware
description tools, of a configurable acceleration architecture based on a
Tensor Core Unit (TCU) composed of multiple Dot Product Units (DPUs). The goal
is to support matrix-oriented computations through different numerical formats
and to enable their comparison in terms of correctness, performance, hardware
cost, numerical behavior, and resilience to faults.

\section{Thesis Objectives}

The main objective of this thesis is to investigate the role of different
numerical formats in hardware accelerators for matrix-oriented computations.
The work aims to study how conventional and emerging number formats can be
supported at hardware level and how their use affects the design of
accelerator architectures.

More specifically, the thesis aims to define a configurable acceleration
architecture for multiply-accumulate operations, integrate it into
representative GPU-like and processor-based environments, and validate its
behavior through simulation and comparison with reference results.


\section{Main Contributions}

Following the mentioned objectives, the main contributions of this thesis can be summarized as follows:

\begin{itemize}

\item A configurable RTL-level Dot Product Unit architecture for
matrix-oriented multiply-accumulate operations under different numerical
formats.

\item A parametric Tensor Core Unit architecture composed of multiple Dot
Product Units and the required control, operand handling, and result
management logic.

\item The integration of the proposed Tensor Core Unit into two different
hardware contexts: the FlexGrip Plus GPU-like architecture and the
PULPino-Klessydra RISC-V core.

\item A functional validation flow based on simulation and comparison with
golden/reference results for verifying the correct behavior of the integrated
accelerator.

\item A hardware-oriented framework that provides a basis for the study of different
numerical formats in terms of correctness, performance, hardware cost,
numerical behavior, and resilience to faults.

\end{itemize}



\begin{itemize}
\item what is the addressed scenario and the addressed problem?
\item why the problem you addressed is interesting and/or practical relevant?

\end{itemize}
Why Number Formats Matter in Hardware Accelerators 


\section{Main Contributions}
\section{Thesis Organization}
